\subsection{Matrix-Induced Maps and Rank}
Throughout this subsection, all matrices are over the same fixed field $\mathbb{F}$, and every matrix size is a positive integer.

\begin{corollary} \label{cor:matrix-rank-nullity} \hlt{Matrix Rank-Nullity}\\
Let $A\in\operatorname{Mat}(m,n,\mathbb{F})$. Then:
\begin{enumerate}[label=(\roman*)]
\setlength{\itemsep}{0pt}
\item $\dim(\ker(A))+\dim(\operatorname{col}(A))=n$; equivalently, $\operatorname{nullity}(A)+\operatorname{rank}(A)=n$.
\item If $m=n$, then $\ker(A)=\{\mathbf{0}\}$ iff $\operatorname{col}(A)=\mathbb{F}^n$.
\end{enumerate}
\begin{proof}
\begin{enumerate}[label=(\roman*)]
\setlength{\itemsep}{0pt}
\item[$(i)$] Write $A=[\mathbf{a}_1\ \cdots\ \mathbf{a}_n]$. The map $T_A:\mathbb{F}^n\to\mathbb{F}^m$ is linear by Proposition \ref{prop:matrix-induced-map-linear}.\\
For $\mathbf{x}=(x_1,\ldots,x_n)^t$, $T_A\mathbf{x}=A\mathbf{x}=\sum\limits_{j=1}^{n}x_j\mathbf{a}_j$ by Proposition \ref{prop:matrix-vector-product-column-combination}.\\
Hence $\operatorname{range}(T_A)=\operatorname{col}(A)$, while $\ker(T_A)=\ker(A)$ holds by $T_A\mathbf{x}=A\mathbf{x}$; $\dim(\ker(A))+\dim(\operatorname{col}(A))=n$ by Theorem \ref{thm:rank-nullity}(ii).\\
The equivalent identity is the same equality under $\operatorname{nullity}(A)=\dim(\ker(A))$ and $\operatorname{rank}(A)=\dim(\operatorname{col}(A))$.
\item[$(ii)$] ($\Rightarrow$) Suppose $m=n$ and $\ker(A)=\{\mathbf{0}\}$. Then $\dim(\operatorname{col}(A))=n=\dim(\mathbb{F}^n)$ by part (i).\\
Since $\operatorname{col}(A)\subseteq\mathbb{F}^n$, $\operatorname{col}(A)=\mathbb{F}^n$ by Proposition \ref{prop:dimension-bounds-and-basis-criteria}(ii).
\item[$(ii)$] ($\Leftarrow$) Suppose $m=n$ and $\operatorname{col}(A)=\mathbb{F}^n$. Then $\dim(\ker(A))=0$ by part (i).\\
The zero subspace and $\ker(A)$ have the same dimension, so $\ker(A)=\{\mathbf{0}\}$ by Proposition \ref{prop:dimension-bounds-and-basis-criteria}(ii).
\end{enumerate}
\hlt{Structural Commentary}: Rank-nullity transfers through the induced map; the square case compares full-dimensional subspaces.\\
\hlt{Key Insight}: A matrix inherits rank-nullity from its coordinate-space map.
\end{proof}
\end{corollary}

\begin{theorem} \label{thm:row-rank-equals-column-rank} \hlt{Row Rank Equals Column Rank}\\
For $A\in\operatorname{Mat}(m,n,\mathbb{F})$, $\dim(\operatorname{row}(A))=\dim(\operatorname{col}(A))$. Equivalently, $\operatorname{rank}(A^t)=\operatorname{rank}(A)$.
\begin{proof}
Let $C=[\mathbf{c}_1\ \cdots\ \mathbf{c}_q]\in\operatorname{Mat}(p,q,\mathbb{F})$ be arbitrary, and set $s\defeq\operatorname{rank}(C)$.\\
If $s=0$, then $\dim(\operatorname{col}(C))=0=\dim(\{\mathbf{0}\})$, so $\operatorname{col}(C)=\{\mathbf{0}\}$ by Proposition \ref{prop:dimension-bounds-and-basis-criteria}(ii). Hence $C=0$ and $\operatorname{rank}(C^t)=0$.\\
Suppose $s>0$. A basis sublist $(\mathbf{c}_{j_1},\ldots,\mathbf{c}_{j_s})$ of the columns exists by Theorem \ref{thm:reducing-spanning-lists}; set $B\defeq[\mathbf{c}_{j_1}\ \cdots\ \mathbf{c}_{j_s}]\in\operatorname{Mat}(p,s,\mathbb{F})$.\\
The relation $(\forall\,k\in\{1,\ldots,q\})(\exists!\,\mathbf{x}_k\in\mathbb{F}^s)\,\mathbf{c}_k=B\mathbf{x}_k$ holds by Proposition \ref{prop:basis-unique-coordinate-representation}.\\
Set $X\defeq[\mathbf{x}_1\ \cdots\ \mathbf{x}_q]\in\operatorname{Mat}(s,q,\mathbb{F})$. Then $C=BX$ by Proposition \ref{prop:matrix-product-by-columns}, and $C^t=X^tB^t$ by Proposition \ref{prop:transpose-reverses-products}.\\
Each column of $C^t$ lies in the span of the $s$ columns of $X^t$ by Proposition \ref{prop:matrix-product-by-columns} and Proposition \ref{prop:matrix-vector-product-column-combination}.\\
Hence $\operatorname{col}(C^t)$ has a basis of at most $s$ members by Theorem \ref{thm:reducing-spanning-lists}, and $\operatorname{rank}(C^t)\leq s=\operatorname{rank}(C)$.\\
Thus $\operatorname{rank}(C^t)\leq\operatorname{rank}(C)$ for the arbitrary matrix $C$ above. For $C=A$, $\operatorname{rank}(A^t)\leq\operatorname{rank}(A)$.\\
For $C=A^t$, $\operatorname{rank}(A)\leq\operatorname{rank}(A^t)$ because $(A^t)^t=A$ by Proposition \ref{prop:elementary-transpose-rules}(i).\\
\hlt{Structural Commentary}: A column-basis factorization proves one rank inequality; transposition supplies the reverse inequality.\\
\hlt{Key Insight}: Transposition preserves the number of independent matrix directions.
\end{proof}
\end{theorem}

\begin{perspective} \label{per:row-column-rank-different-hosts} \hlt{Row and Column Spaces Have Different Hosts}\\
For $A\in\operatorname{Mat}(m,n,\mathbb{F})$, $\operatorname{row}(A)=\operatorname{col}(A^t)\subseteq\mathbb{F}^n$ and $\operatorname{col}(A)\subseteq\mathbb{F}^m$.\\
When $m\neq n$, Theorem \ref{thm:row-rank-equals-column-rank} identifies their dimensions, not the subspaces in their different ambient spaces.
\end{perspective}

\begin{theorem} \label{thm:matrix-four-fundamental-subspaces-dimensions} \hlt{Four Fundamental Matrix Subspaces}\\
Let $A\in\operatorname{Mat}(m,n,\mathbb{F})$ and let $r\defeq\operatorname{rank}(A)$. Then:
\begin{enumerate}[label=(\roman*)]
\setlength{\itemsep}{0pt}
\item $\operatorname{row}(A),\ker(A)\subseteq\mathbb{F}^n$ and $\operatorname{col}(A),\ker(A^t)\subseteq\mathbb{F}^m$.
\item $\dim(\operatorname{col}(A))=\dim(\operatorname{row}(A))=r$, $\dim(\ker(A))=n-r$, and $\dim(\ker(A^t))=m-r$.
\end{enumerate}
\begin{proof}
\begin{enumerate}[label=(\roman*)]
\setlength{\itemsep}{0pt}
\item The row and column inclusions hold by their entry lengths, while the kernel inclusions hold by the domains of $T_A$ and $T_{A^t}$.
\item The equality $\dim(\operatorname{col}(A))=r$ holds by the definition of rank, and $\dim(\operatorname{row}(A))=r=\operatorname{rank}(A^t)$ by Theorem \ref{thm:row-rank-equals-column-rank}.\\
Hence $\dim(\ker(A))=n-r$ and $\dim(\ker(A^t))=m-r$ by Corollary \ref{cor:matrix-rank-nullity}(i) for $A$ and $A^t$, respectively.
\end{enumerate}
\hlt{Structural Commentary}: Ambient sizes place the subspaces; row-column equality and rank-nullity determine their dimensions.\\
\hlt{Key Insight}: One rank parameter controls both image and kernel dimensions.
\end{proof}
\end{theorem}

\begin{illustrationfigure}[H]
  \scalebox{0.9}{\begin{tikzpicture}[
  scale=0.90,
  transform shape,
  >=Stealth,
  every node/.style={font=\small},
  host/.style={draw=gray!50, fill=gray!3, rounded corners=3pt, line width=0.55pt,
    minimum width=4.05cm, minimum height=2.05cm},
  entry/.style={draw=gray!60, rounded corners=2pt, line width=0.5pt,
    minimum width=1.72cm, minimum height=0.78cm, align=center},
  active/.style={entry, fill=teal!8},
  null/.style={entry, fill=blue!7},
  map/.style={->, gray!65!black, line width=0.75pt}
]
  \node[host] (Fn) at (0,0) {};
  \node[anchor=north west] at (-1.82,0.88) {$\mathbb{F}^n$};
  \node[host] (Fm) at (5.6,0) {};
  \node[anchor=north west] at (3.78,0.88) {$\mathbb{F}^m$};

  \node[null] at (-0.98,-0.05) {$\ker(A)$\\[-1pt]{\footnotesize $n-r$}};
  \node[active] at (0.98,-0.05) {$\operatorname{row}(A)$\\[-1pt]{\footnotesize $r$}};
  \node[active] at (4.62,-0.05) {$\operatorname{col}(A)$\\[-1pt]{\footnotesize $r$}};
  \node[null] at (6.58,-0.05) {$\ker(A^t)$\\[-1pt]{\footnotesize $m-r$}};

  \draw[map] (2.1,0.4) -- (3.5,0.4) node[midway, above] {$A$};
  \draw[map] (3.5,-0.4) -- (2.1,-0.4) node[midway, below] {$A^t$};
\end{tikzpicture}
}
  \caption{Hosts and dimensions of the four fundamental subspaces, with arrows for $A:\mathbb{F}^n\to\mathbb{F}^m$ and $A^t:\mathbb{F}^m\to\mathbb{F}^n$.}
  \label{fig:four-fundamental-subspaces-hosts}
\end{illustrationfigure}

Figure \ref{fig:four-fundamental-subspaces-hosts} illustrates Theorem \ref{thm:matrix-four-fundamental-subspaces-dimensions}. Its boxes encode only ambient hosts and dimensions, not direct sums or orthogonality.

\begin{corollary} \label{cor:homogeneous-system-more-variables-than-equations} \hlt{Homogeneous Systems with More Variables Than Equations}\\
If $A\in\operatorname{Mat}(m,n,\mathbb{F})$ with $n>m$, then $(\exists\,\mathbf{x}\in\mathbb{F}^n\setminus\{\mathbf{0}\})\,A\mathbf{x}=\mathbf{0}$.
\begin{proof}
The map $T_A:\mathbb{F}^n\to\mathbb{F}^m$ is linear by Proposition \ref{prop:matrix-induced-map-linear} and is not injective by Corollary \ref{cor:map-to-smaller-space-not-injective}.\\
Thus $\ker(T_A)\neq\{\mathbf{0}\}$ by Proposition \ref{prop:injective-iff-kernel-zero}. Choose $\mathbf{x}\in\ker(T_A)\setminus\{\mathbf{0}\}$. Then $A\mathbf{x}=T_A\mathbf{x}=\mathbf{0}$.\\
\hlt{Structural Commentary}: The dimension obstruction transfers from the induced map to the homogeneous system.\\
\hlt{Key Insight}: More variables force a nontrivial nullspace.
\end{proof}
\end{corollary}

\begin{corollary} \label{cor:inhomogeneous-system-more-equations-than-variables} \hlt{Inhomogeneous Systems with More Equations Than Variables}\\
If $A\in\operatorname{Mat}(m,n,\mathbb{F})$ with $m>n$, then $(\exists\,\mathbf{b}\in\mathbb{F}^m)(\forall\,\mathbf{x}\in\mathbb{F}^n)\,A\mathbf{x}\neq\mathbf{b}$.
\begin{proof}
The map $T_A:\mathbb{F}^n\to\mathbb{F}^m$ is linear by Proposition \ref{prop:matrix-induced-map-linear} and is not surjective by Corollary \ref{cor:map-to-larger-space-not-surjective}.\\
Choose $\mathbf{b}\in\mathbb{F}^m\setminus\operatorname{range}(T_A)$. Then $(\forall\,\mathbf{x}\in\mathbb{F}^n)\,A\mathbf{x}=T_A\mathbf{x}\neq\mathbf{b}$.\\
\hlt{Structural Commentary}: The dimension obstruction leaves a right-hand side outside the induced map's range.\\
\hlt{Key Insight}: More equations force an unreachable right-hand side.
\end{proof}
\end{corollary}
